% 六封邮件及同一星标保留要求；删除/归档是互斥方案，不是先后步骤。
\begin{tikzpicture}[figure picture,foundation picture,
 trust node/.style={draw=FigureStroke,line width=.7pt,rounded corners=1.5pt,align=center,minimum height=13mm,inner sep=2mm}]
 \path[use as bounding box] (0,0) rectangle (163.2,88);
 \node at (81.6,82) {星标内容应保留};
 \draw[draw=FigureBlue,line width=.7pt,rounded corners=1.5pt] (63,55) rectangle(100,77);
 \foreach \x in {66,78,90}{\foreach \y in {59,68}{
  \draw[draw=FigureStroke,fill=white,line width=.65pt] (\x,\y) rectangle({\x+8},{\y+5});
  \draw[draw=FigureStroke,line width=.65pt] (\x,{\y+5})--({\x+4},{\y+1.5})--({\x+8},{\y+5});
 }}
 \node[text=LancetYellow,inner sep=0pt] at (73.5,73) {$\star$};
 \node[trust node,fill=FigureRedFill,minimum width=49mm] (delete) at (37,39) {删除全部六封邮件};
 \node[trust node,fill=MutedGreenFill,minimum width=49mm] (archive) at (126,39) {归档全部六封邮件};
 \node[trust node,fill=FigureRedFill,minimum width=62mm,minimum height=19mm] (lost) at (37,12) {收件箱为空\\星标内容丢失};
 \node[trust node,fill=MutedGreenFill,minimum width=62mm,minimum height=19mm] (kept) at (126,12) {收件箱为空\\星标内容仍可检索};
 \draw[foundation flow] (81.6,55)--(81.6,50)-|(delete.north);
 \draw[foundation flow] (81.6,50)-|(archive.north);
 \draw[foundation flow] (delete.south)--(lost.north);
 \draw[foundation flow] (archive.south)--(kept.north);
\end{tikzpicture}
